% Options for packages loaded elsewhere
\PassOptionsToPackage{unicode}{hyperref}
\PassOptionsToPackage{hyphens}{url}
\documentclass[
]{article}
\usepackage{xcolor}
\usepackage[margin=1in]{geometry}
\usepackage{amsmath,amssymb}
\setcounter{secnumdepth}{-\maxdimen} % remove section numbering
\usepackage{iftex}
\ifPDFTeX
  \usepackage[T1]{fontenc}
  \usepackage[utf8]{inputenc}
  \usepackage{textcomp} % provide euro and other symbols
\else % if luatex or xetex
  \usepackage{unicode-math} % this also loads fontspec
  \defaultfontfeatures{Scale=MatchLowercase}
  \defaultfontfeatures[\rmfamily]{Ligatures=TeX,Scale=1}
\fi
\usepackage{lmodern}
\ifPDFTeX\else
  % xetex/luatex font selection
\fi
% Use upquote if available, for straight quotes in verbatim environments
\IfFileExists{upquote.sty}{\usepackage{upquote}}{}
\IfFileExists{microtype.sty}{% use microtype if available
  \usepackage[]{microtype}
  \UseMicrotypeSet[protrusion]{basicmath} % disable protrusion for tt fonts
}{}
\makeatletter
\@ifundefined{KOMAClassName}{% if non-KOMA class
  \IfFileExists{parskip.sty}{%
    \usepackage{parskip}
  }{% else
    \setlength{\parindent}{0pt}
    \setlength{\parskip}{6pt plus 2pt minus 1pt}}
}{% if KOMA class
  \KOMAoptions{parskip=half}}
\makeatother
\usepackage{color}
\usepackage{fancyvrb}
\newcommand{\VerbBar}{|}
\newcommand{\VERB}{\Verb[commandchars=\\\{\}]}
\DefineVerbatimEnvironment{Highlighting}{Verbatim}{commandchars=\\\{\}}
% Add ',fontsize=\small' for more characters per line
\usepackage{framed}
\definecolor{shadecolor}{RGB}{248,248,248}
\newenvironment{Shaded}{\begin{snugshade}}{\end{snugshade}}
\newcommand{\AlertTok}[1]{\textcolor[rgb]{0.94,0.16,0.16}{#1}}
\newcommand{\AnnotationTok}[1]{\textcolor[rgb]{0.56,0.35,0.01}{\textbf{\textit{#1}}}}
\newcommand{\AttributeTok}[1]{\textcolor[rgb]{0.13,0.29,0.53}{#1}}
\newcommand{\BaseNTok}[1]{\textcolor[rgb]{0.00,0.00,0.81}{#1}}
\newcommand{\BuiltInTok}[1]{#1}
\newcommand{\CharTok}[1]{\textcolor[rgb]{0.31,0.60,0.02}{#1}}
\newcommand{\CommentTok}[1]{\textcolor[rgb]{0.56,0.35,0.01}{\textit{#1}}}
\newcommand{\CommentVarTok}[1]{\textcolor[rgb]{0.56,0.35,0.01}{\textbf{\textit{#1}}}}
\newcommand{\ConstantTok}[1]{\textcolor[rgb]{0.56,0.35,0.01}{#1}}
\newcommand{\ControlFlowTok}[1]{\textcolor[rgb]{0.13,0.29,0.53}{\textbf{#1}}}
\newcommand{\DataTypeTok}[1]{\textcolor[rgb]{0.13,0.29,0.53}{#1}}
\newcommand{\DecValTok}[1]{\textcolor[rgb]{0.00,0.00,0.81}{#1}}
\newcommand{\DocumentationTok}[1]{\textcolor[rgb]{0.56,0.35,0.01}{\textbf{\textit{#1}}}}
\newcommand{\ErrorTok}[1]{\textcolor[rgb]{0.64,0.00,0.00}{\textbf{#1}}}
\newcommand{\ExtensionTok}[1]{#1}
\newcommand{\FloatTok}[1]{\textcolor[rgb]{0.00,0.00,0.81}{#1}}
\newcommand{\FunctionTok}[1]{\textcolor[rgb]{0.13,0.29,0.53}{\textbf{#1}}}
\newcommand{\ImportTok}[1]{#1}
\newcommand{\InformationTok}[1]{\textcolor[rgb]{0.56,0.35,0.01}{\textbf{\textit{#1}}}}
\newcommand{\KeywordTok}[1]{\textcolor[rgb]{0.13,0.29,0.53}{\textbf{#1}}}
\newcommand{\NormalTok}[1]{#1}
\newcommand{\OperatorTok}[1]{\textcolor[rgb]{0.81,0.36,0.00}{\textbf{#1}}}
\newcommand{\OtherTok}[1]{\textcolor[rgb]{0.56,0.35,0.01}{#1}}
\newcommand{\PreprocessorTok}[1]{\textcolor[rgb]{0.56,0.35,0.01}{\textit{#1}}}
\newcommand{\RegionMarkerTok}[1]{#1}
\newcommand{\SpecialCharTok}[1]{\textcolor[rgb]{0.81,0.36,0.00}{\textbf{#1}}}
\newcommand{\SpecialStringTok}[1]{\textcolor[rgb]{0.31,0.60,0.02}{#1}}
\newcommand{\StringTok}[1]{\textcolor[rgb]{0.31,0.60,0.02}{#1}}
\newcommand{\VariableTok}[1]{\textcolor[rgb]{0.00,0.00,0.00}{#1}}
\newcommand{\VerbatimStringTok}[1]{\textcolor[rgb]{0.31,0.60,0.02}{#1}}
\newcommand{\WarningTok}[1]{\textcolor[rgb]{0.56,0.35,0.01}{\textbf{\textit{#1}}}}
\usepackage{graphicx}
\makeatletter
\newsavebox\pandoc@box
\newcommand*\pandocbounded[1]{% scales image to fit in text height/width
  \sbox\pandoc@box{#1}%
  \Gscale@div\@tempa{\textheight}{\dimexpr\ht\pandoc@box+\dp\pandoc@box\relax}%
  \Gscale@div\@tempb{\linewidth}{\wd\pandoc@box}%
  \ifdim\@tempb\p@<\@tempa\p@\let\@tempa\@tempb\fi% select the smaller of both
  \ifdim\@tempa\p@<\p@\scalebox{\@tempa}{\usebox\pandoc@box}%
  \else\usebox{\pandoc@box}%
  \fi%
}
% Set default figure placement to htbp
\def\fps@figure{htbp}
\makeatother
\setlength{\emergencystretch}{3em} % prevent overfull lines
\providecommand{\tightlist}{%
  \setlength{\itemsep}{0pt}\setlength{\parskip}{0pt}}
\usepackage{bookmark}
\IfFileExists{xurl.sty}{\usepackage{xurl}}{} % add URL line breaks if available
\urlstyle{same}
\hypersetup{
  pdftitle={Assignment 6 Report},
  pdfauthor={Taylor Falk},
  hidelinks,
  pdfcreator={LaTeX via pandoc}}

\title{Assignment 6 Report}
\author{Taylor Falk}
\date{1/24/2022}

\begin{document}
\maketitle

\section{Comparison of Three Differential Expression
Packages}\label{comparison-of-three-differential-expression-packages}

We are interested in examining how three differential expression
packages (DESeq2, edgeR, and Limma) operate within R. Complete the code
blocks below and use this document as a reference to recreate the
appropriate results. This report document should be completed after you
fill in the functions in \texttt{main.R} and you satisfy the tests in
\texttt{test\_main.R} to the best of your ability. Once this document is
completed, \texttt{Knit} it and push your repository to your student
GitHub repository.

This assignment will involve a decent amount of documentation reading,
but the relevant information is linked in the \texttt{main.R} function
descriptions.

\textbf{WARNING} For SCC Users

One of the libraries we utilize in this project, ggVennDiagram, requires
some packages that are not enabled by default on SCC. If you install
this package while using RStudio on your machine, you are likely not to
run into this problem. In order to load these packages, we will enter
two packages into the RStudio Server prompt when starting it on SCC (big
shout outs to Mae Rose for solving this one!).

Open the RStudio Server launcher on the SCC for web interface, and click
on ``Select Modules''.
\pandocbounded{\includegraphics[keepaspectratio]{mkd_img/scc1.png}}

While in this submenu, search and select the modules \texttt{proj/8.2.1}
and \texttt{sqlite3/3.37.2}. Your screen should look like the below
image.
\pandocbounded{\includegraphics[keepaspectratio]{mkd_img/scc2.png}}

Finally, close out, your page should look like this final image. Proceed
to start the session as normal.
\pandocbounded{\includegraphics[keepaspectratio]{mkd_img/scc3.png}}

If you like aesthetics, feel free to remove the above warning section
from your final report.

\subsection{Setup}\label{setup}

\begin{Shaded}
\begin{Highlighting}[]
\FunctionTok{source}\NormalTok{(}\StringTok{"main.R"}\NormalTok{) }\CommentTok{\# load all our wonderful functions}
\NormalTok{knitr}\SpecialCharTok{::}\NormalTok{opts\_chunk}\SpecialCharTok{$}\FunctionTok{set}\NormalTok{(}\AttributeTok{echo =} \ConstantTok{FALSE}\NormalTok{) }\CommentTok{\# no markdown for you \textless{}3}
\end{Highlighting}
\end{Shaded}

Load the data using \texttt{load\_n\_trim()} and store into a variable.

\begin{verbatim}
##                      vP0_1 vP0_2 vAd_1 vAd_2
## ENSMUSG00000102693.2     0     0     0     0
## ENSMUSG00000064842.3     0     0     0     0
## ENSMUSG00000051951.6    19    24     7     5
## ENSMUSG00000102851.2     0     0     0     0
## ENSMUSG00000103377.2     1     0     1     0
## ENSMUSG00000104017.2     0     3     0     0
\end{verbatim}

\subsection{Run three packages}\label{run-three-packages}

In the following three sections, use your functions to complete the
three differential expression analyses for DESeq2, edgeR, and
Limma-voom.

\subsubsection{DESeq2}\label{deseq2}

We include the output of \texttt{run\_deseq()} as well as the contents
of the \texttt{coldata} input. The \texttt{colData} parameter for DESeq2
describes the experimental setup and which samples belong to which part
of the experiment. We will be using the recommended count filter of 10
(from the DESeq2 documentation). A sample of the DESeq2 output is
included, but note that the order of your genes and the values in the
columns may differ slightly, please ensure the data is the right shape
and has columns labeled correctly. A warning from DESeq2 is normal.
Store DESeq2 results into a variable.

\begin{verbatim}
## coldata:
\end{verbatim}

\begin{verbatim}
##       condition
## vP0_1      day0
## vP0_2      day0
## vAd_1     adult
## vAd_2     adult
\end{verbatim}

\begin{verbatim}
## Warning in DESeqDataSet(se, design = design, ignoreRank): some variables in
## design formula are characters, converting to factors
\end{verbatim}

\begin{verbatim}
## estimating size factors
\end{verbatim}

\begin{verbatim}
## estimating dispersions
\end{verbatim}

\begin{verbatim}
## gene-wise dispersion estimates
\end{verbatim}

\begin{verbatim}
## mean-dispersion relationship
\end{verbatim}

\begin{verbatim}
## final dispersion estimates
\end{verbatim}

\begin{verbatim}
## fitting model and testing
\end{verbatim}

\begin{verbatim}
## deseq results:
\end{verbatim}

\begin{verbatim}
## log2 fold change (MLE): condition day0 vs adult 
## Wald test p-value: condition day0 vs adult 
## DataFrame with 6 rows and 6 columns
##                        baseMean log2FoldChange     lfcSE      stat      pvalue
##                       <numeric>      <numeric> <numeric> <numeric>   <numeric>
## ENSMUSG00000051951.6   12.45040       1.350880  0.908351  1.487178 1.36968e-01
## ENSMUSG00000102331.2    4.73568       0.395338  1.389527  0.284513 7.76017e-01
## ENSMUSG00000025900.14  22.39009      -3.833022  0.833865 -4.596695 4.29245e-06
## ENSMUSG00000025902.14 174.44409      -0.247347  0.292435 -0.845820 3.97653e-01
## ENSMUSG00000104238.2    2.56484      -1.134668  2.097294 -0.541015 5.88497e-01
## ENSMUSG00000102269.2    7.87344       0.576799  1.156710  0.498655 6.18023e-01
##                              padj
##                         <numeric>
## ENSMUSG00000051951.6  2.63320e-01
## ENSMUSG00000102331.2  8.73332e-01
## ENSMUSG00000025900.14 2.75823e-05
## ENSMUSG00000025902.14 5.73336e-01
## ENSMUSG00000104238.2           NA
## ENSMUSG00000102269.2  7.63028e-01
\end{verbatim}

\subsubsection{EdgeR}\label{edger}

EdgeR's method of describing an experiment is more simple than DESeq2.
The default results for EdgeR \textbf{do not} include p-adjusted values,
but we can add them ourselves. Once we have the results stored, we can
add a new column by using the \texttt{stats::p.adjust()} function with
\texttt{method\ =\ \textquotesingle{}BH\textquotesingle{}}.

\begin{verbatim}
## group:
\end{verbatim}

\begin{verbatim}
## 1 1 2 2
\end{verbatim}

\begin{verbatim}
## Using classic mode.
\end{verbatim}

\begin{verbatim}
## edger without padj:
\end{verbatim}

\begin{verbatim}
##                            logFC     logCPM       PValue
## ENSMUSG00000051951.6  -1.3336045 -0.1453460 1.532615e-02
## ENSMUSG00000025900.14  3.7815113  0.3977791 1.520214e-12
## ENSMUSG00000025902.14  0.2514873  3.3930556 3.215474e-01
\end{verbatim}

\begin{verbatim}
## edger _with_ padj:
\end{verbatim}

\begin{verbatim}
##                            logFC     logCPM       PValue         padj
## ENSMUSG00000051951.6  -1.3336045 -0.1453460 1.532615e-02 3.149080e-02
## ENSMUSG00000025900.14  3.7815113  0.3977791 1.520214e-12 1.435984e-11
## ENSMUSG00000025902.14  0.2514873  3.3930556 3.215474e-01 4.235719e-01
\end{verbatim}

\subsubsection{Limma with voom}\label{limma-with-voom}

Finally, we create an experimental design variable for Limma. We
\emph{will} be using voom in this analysis, so set your voom argument to
\texttt{TRUE}. If your limma/voom functions include a plot, that's fine!
Bonus plots are always fun for cramming into your supplemental
materials.

\begin{verbatim}
## design:
\end{verbatim}

\begin{verbatim}
##       day0 day0vsadult
## vP0_1    1           0
## vP0_2    1           0
## vAd_1    1           1
## vAd_2    1           1
\end{verbatim}

\begin{verbatim}
## group:
\end{verbatim}

\begin{verbatim}
## [1] 1 1 2 2
## Levels: 1 2
\end{verbatim}

\begin{verbatim}
## Limma results:
\end{verbatim}

\begin{verbatim}
##                           logFC  AveExpr         t      P.Value    adj.P.Val
## ENSMUSG00000026418.17 -9.817581 5.534673 -39.42302 6.856393e-09 3.841258e-05
## ENSMUSG00000002500.16  6.177204 5.616861  37.76303 9.023170e-09 3.841258e-05
## ENSMUSG00000051855.16 -5.358358 5.891389 -36.10231 1.202260e-08 3.841258e-05
## ENSMUSG00000042828.13  5.322768 6.682673  35.42601 1.356413e-08 3.841258e-05
## ENSMUSG00000021622.4   4.551307 9.187568  34.73656 1.537570e-08 3.841258e-05
## ENSMUSG00000031972.6   4.407266 8.129513  34.09673 1.731116e-08 3.841258e-05
##                               B
## ENSMUSG00000026418.17  8.655442
## ENSMUSG00000002500.16 10.320906
## ENSMUSG00000051855.16 10.240130
## ENSMUSG00000042828.13 10.566226
## ENSMUSG00000021622.4  10.628166
## ENSMUSG00000031972.6  10.508273
\end{verbatim}

\subsection{Plots}\label{plots}

We don't want to worry our little heads plotting \emph{all} of the data,
so we will just trim it down to the 1,000 most significant p-values.
Order each of the three results set by \textbf{p-value} and take the top
(smallest) 1,000 rows. Note that we are \emph{plotting} the
\textbf{p-adjusted} values but we are \emph{sorting} the
\textbf{p-values}. This is an important distinction.

\begin{verbatim}
## Limma shape before:
\end{verbatim}

\begin{verbatim}
## 'data.frame':    15548 obs. of  6 variables:
##  $ logFC    : num  -9.82 6.18 -5.36 5.32 4.55 ...
##  $ AveExpr  : num  5.53 5.62 5.89 6.68 9.19 ...
##  $ t        : num  -39.4 37.8 -36.1 35.4 34.7 ...
##  $ P.Value  : num  6.86e-09 9.02e-09 1.20e-08 1.36e-08 1.54e-08 ...
##  $ adj.P.Val: num  3.84e-05 3.84e-05 3.84e-05 3.84e-05 3.84e-05 ...
##  $ B        : num  8.66 10.32 10.24 10.57 10.63 ...
## NULL
\end{verbatim}

\begin{verbatim}
## Limma shape after:
\end{verbatim}

\begin{verbatim}
## 'data.frame':    1000 obs. of  6 variables:
##  $ logFC    : num  -9.82 6.18 -5.36 5.32 4.55 ...
##  $ AveExpr  : num  5.53 5.62 5.89 6.68 9.19 ...
##  $ t        : num  -39.4 37.8 -36.1 35.4 34.7 ...
##  $ P.Value  : num  6.86e-09 9.02e-09 1.20e-08 1.36e-08 1.54e-08 ...
##  $ adj.P.Val: num  3.84e-05 3.84e-05 3.84e-05 3.84e-05 3.84e-05 ...
##  $ B        : num  8.66 10.32 10.24 10.57 10.63 ...
## NULL
\end{verbatim}

\begin{verbatim}
## head(deseq_results), notice the order of the p-values nad p-adjusted values:
\end{verbatim}

We want to create a venn diagram using the \texttt{ggVennDiagram}
package. This package asks for a list object as input, so we can create
a list that contains the three named collections of genes from our
separate packages. Calling row.names() will give us the names of all the
genes. More info here:
\href{https://r-charts.com/part-whole/ggvenndiagram/}{ggVennDiagram}

\begin{figure}

{\centering \includegraphics{report_files/figure-latex/Venn diagram-1} 

}

\caption{Figure 1 - A venn diagram comparing the top 1,000 differentially expressed genes from three different R packages ordered by ascending p-value. A Lighter shade of blue indicates a higher proportion of shared genes.}\label{fig:Venn diagram}
\end{figure}

While plots can be used to create publication ready images, we can also
utilize them to check our data and confirm our assumptions. Limma seems
to have a number of selected genes that don't match DESeq2 and edgeR, so
we're interested to see how the distributions of p-values compare to one
another. We did not write this plot as a function, but can use the data
we generate with \texttt{combine\_pval()} to \texttt{facet\_wrap()} the
three data sets together. Use facet wrap with the
\texttt{\textasciitilde{}\ package} parameter to recreate the figure
below.

\begin{figure}

{\centering \includegraphics{report_files/figure-latex/Facet plot-1} 

}

\caption{Figure 2 - Three histogram plots comparing the distribution of p-values for three differential expression packages. Note that edgeR and DESeq2 have their values concentrated closer to 0 than the Limma package results.}\label{fig:Facet plot}
\end{figure}

Use \texttt{create\_facets()} and \texttt{theme\_plot()} to create the
final plot.

\begin{figure}

{\centering \includegraphics{report_files/figure-latex/Final plot-1} 

}

\caption{Figure 3 - A volcano plot of the top 1,000 genes (ranked by p-value) comparing the log<sub>2</sub> fold-change with the adjusted p-value. Adjusted p-values above 1e-100 are highlighted in red.}\label{fig:Final plot}
\end{figure}

\textbf{Bonus} Compare Limma Without Voom

Completing the above cell is enough to finish the assignment, but there
is one element of the documentation for Limma that you may be interested
in looking at (entirely optionally, of course). \texttt{voom}, the
component of Limma that we utilize, is not exactly \emph{necessary}. The
documentation even suggests: \textgreater{} If the sequencing depth is
reasonably consistent across the RNA samples, then the simplest and most
robust approach to differential exis to use limma-trend.

Is this appropriate for our data? How can you tell? If you're
interested, design a function to run limma-trend \emph{instead} of voom,
and compare using the plots above. What does using voom do to the data?
Again, this is a bonus section and {\textbf{is not}} necessary to
complete assignment 6. Feel free to create the venn diagram, or other
plots used in this assignment.\\

\begin{figure}

{\centering \includegraphics{report_files/figure-latex/Voomless bonus-1} 

}

\caption{Figure 4+ - A venn diagram comparing the top 1,000 differentially expressed genes from four different R packages ordered by ascending p-value.}\label{fig:Voomless bonus}
\end{figure}

\end{document}
